\renewcommand{\thetable}{3.A.1}
{\LFXcuerpotabla\setlength{\LTpre}{12pt}\setlength{\LTpost}{16pt}%
\begin{longtable}{@{}F{0.062} F{0.086} F{0.232} F{0.189} F{0.228} F{0.098} F{0.105}@{}}
\caption{Requerimientos funcionales del caso. Fuente: elaboración propia a partir de los requerimientos de Distribuidora Puelche y del Anexo 2.1 del Subdocumento 2 (LafroX, 2026a).}\label{tab:3-A-1}\\[2pt]
\toprule \cab{ID} & \cab{Actor} & \cab{Descripción} & \cab{Precondición} & \cab{Resultado esperado} & \cab{Prioridad / etapa} & \cab{Origen} \\ \midrule
\endfirsthead
\multicolumn{7}{@{}l@{}}{\fontsize{9pt}{10.5pt}\selectfont\color{lafroxGris}Tabla \thetable\ --- continuación}\\[3pt]
\toprule \cab{ID} & \cab{Actor} & \cab{Descripción} & \cab{Precondición} & \cab{Resultado esperado} & \cab{Prioridad / etapa} & \cab{Origen} \\ \midrule
\endhead
\midrule
\multicolumn{7}{@{}l@{}}{\fontsize{9pt}{10.5pt}\selectfont\color{lafroxGris}continúa en la página siguiente}\\
\endfoot
\bottomrule
\endlastfoot
  RF-01.01 & Jefa de Bodega & El sistema debe registrar la recepción física de un proveedor validando las cantidades recibidas por SKU contra la OC y generando una alerta por cada discrepancia, sin bloquear el cierre. & Existe una OC emitida y confirmada al proveedor en el ERP con sus líneas de SKU, y el recepcionista inicia el registro de la recepción física del proveedor en el andén (RF-01.02). & El comparativo pedido vs. recibido queda registrado al cerrar la recepción y toda discrepancia queda disponible en el módulo de calidad. & Crítica / E1 &  \\
  RF-01.02 & Jefa de Bodega & El sistema debe exigir el número de lote del proveedor en campo estructurado para todo SKU con trazabilidad obligatoria, capturándolo por escaneo GS1 (AI 10) o ingreso manual, y bloqueando el avance si el campo está vacío. & El SKU está definido en la ficha de producto con el indicador de trazabilidad obligatoria (RF-01.10) y el recepcionista se encuentra registrando la recepción del lote (RF-01.01). & El 100 \% de los ítems con trazabilidad obligatoria registran lote; campo vacío muestra "Lote requerido para este producto" e impide continuar. & Crítica / E1 & . 7.3, \\
  RF-01.03 & Jefa de Bodega & El sistema debe registrar la fecha de vencimiento en formato DD/MM/AAAA para todo SKU con control de vencimiento habilitado, vinculándola al lote y a la ubicación asignada como entrada para la lógica FEFO. & El SKU tiene control de vencimiento habilitado en su ficha de producto (RF-01.10) y la recepción está en curso con el lote identificado (RF-01.02). & El 100 \% de los SKU con control de vencimiento habilitado tienen fecha registrada y asociada al lote al cerrar la recepción; campo vacío bloquea el avance. & Crítica / E1 & . 4.8 \\
  RF-01.04 & Jefa de Bodega & El sistema debe permitir registrar una no conformidad durante la recepción seleccionando su tipo desde una lista predefinida e ingresando SKU, lote, cantidad afectada y al menos una fotografía capturada desde el dispositivo. & El recepcionista detecta una discrepancia o anomalía de calidad durante la recepción (RF-01.01) y el SKU/lote afectado está identificado en la recepción en curso. & La no conformidad queda registrada con todos los campos requeridos y la fotografía es descargable desde el módulo de calidad. & Alta / E1 & . 4.9, Cap 8 (Entrevista de Calidad), RT-17.06 \\
  RF-01.05 & Jefa de Bodega & Al confirmar una no conformidad, el sistema debe cambiar automáticamente a estado ‘Cuarentena’ la cantidad afectada del SKU/lote correspondiente, bloqueándola para picking, y notificar al perfil de calidad mediante una alerta en la aplicación. & Existe una no conformidad confirmada con SKU, lote y cantidad afectada definidos (RF-01.04). & La cantidad del lote queda bloqueado para picking de forma inmediata y el perfil de calidad recibe la notificación en $\le$ 5 min & Alta / E1 & . 4.9, RT-17.06 \\
  RF-01.06 & Jefa de Bodega & El sistema debe decodificar códigos de barras GS1 interpretando los AIs 01 (GTIN), 10 (lote), 17 (vencimiento) y 00 (SSCC), autocompletando los campos en pantalla; ante código inválido o no registrado, habilita el ingreso manual. & El dispositivo de recepción cuenta con cámara o lector de código de barras operativo (RT-17.06) y el rotulado GS1 del producto o unidad logística está disponible para escaneo. & Código GS1 válido autocompleta los campos en pantalla; código no encontrado muestra "Producto no encontrado" y habilita ingreso manual. & Alta / E1 & Caso, RT-17.06, Cap. 16.2 (GS1) \\
  RF-01.07 & Jefa de Bodega & Al finalizar el registro de un pallet, el sistema debe generar automáticamente un SSCC bajo estándar GS1-128 vinculado al GTIN, lote, fecha de vencimiento y fecha de recepción, e imprimir la etiqueta en la impresora térmica del andén. & El pallet quedó registrado con GTIN, lote y fecha de vencimiento (RF-01.02, RF-01.03, RF-01.06) y existe una impresora térmica del andén operativa y conectada. & El SSCC se genera y la etiqueta se imprime con los campos GTIN (AI 01), lote (AI 10), fecha de vencimiento (AI 17) y fecha de recepción (AI 11). & Media / E1 & ), RT-05.23 \\
  RF-01.08 & Jefa de Bodega & El sistema debe permitir asociar un SSCC a una dirección de almacenamiento (pasillo-columna-nivel) escaneando ambas etiquetas y validando la compatibilidad térmica antes de confirmar. & Existe un SSCC generado (RF-01.07), la topología del CD está configurada con zonas térmicas y posiciones (RF-02.01) y el operador dispone de las etiquetas de SSCC y ubicación para escanear. & La asociación SSCC–ubicación queda registrada; intento de ubicación incompatible es bloqueado con mensaje descriptivo. & Media / E1 & . 12 (GS1), RT-05.23 \\
  RF-01.09 & Jefa de Bodega & Al cerrar una recepción, el sistema debe enviar automáticamente la transacción de ingreso de inventario al ERP mediante su interfaz de integración; si el ERP no está disponible, encola la transacción para reintento automático. & La recepción fue cerrada con su comparativo pedido vs. recibido (RF-01.01) y la interfaz de integración con el ERP está configurada y disponible. & La conciliación diaria no muestra diferencias entre el inventario del sistema y el del ERP; ante ERP no disponible, el sistema informa al usuario y reintenta sin intervención. & Crítica / E1 &  \\
  RF-01.10 & Gerente Comercial / Gerente de Operaciones & El sistema debe permitir crear y mantener fichas de producto con atributos obligatorios: descripción, formato, GTIN, SKU interno, tipo de almacenamiento, temperatura mínima y máxima (°C), peso neto (kg), dimensiones (cm), días de vida útil, indicadores de trazabilidad y control de vencimiento (sí/no), clasificación sanitaria y RUT del proveedor. & El administrador de catálogo / jefe de abastecimiento tiene perfil con permisos de creación y mantención de fichas de producto. & El sistema no permite guardar una ficha con campos obligatorios vacíos y registra historial de cambios de GTIN o SKU con fecha y usuario. & Crítica / E1 & Caso, Cap 14.1 (pasa de 8.400 a 9.500 SKu), Cap 8 (Hugo), Cap 16.1 (N11) \\
  RF-01.11 & Jefa de Bodega & El sistema debe validar la compatibilidad térmica entre el tipo de almacenamiento del SKU y la zona de temperatura de la ubicación destino al realizar un putaway, bloqueando la operación si no son compatibles. & El SKU tiene tipo de almacenamiento definido en su ficha de producto (RF-01.10), la ubicación destino tiene zona térmica configurada (RF-02.01) y el operador ejecuta un putaway con el SSCC o unidad a ubicar (RF-01.08). & Intento de ubicar un SKU congelado fuera de zona congelada es bloqueado con mensaje "Incompatibilidad térmica: requiere zona congelado". & Crítica / E1 & . 4.9, Cap. 10 (restricción 2) \\
  RF-02.01 & Jefe de TI / Jefa de Bodega & El sistema debe permitir registrar y configurar cada instalación de la red logística (centros de distribución y plataformas de cross-docking), definiendo su topología interna: zonas operativas (seco con rango 15–25\,°C, refrigerado con rango 2–8\,°C, congelado con rango $-$18 a $-$25\,°C), pasillos, racks, niveles y posiciones individuales de almacenamiento. La configuración debe estar disponible en modo offline. & El administrador tiene perfil con permisos de configuración de instalaciones. & (a) Al configurar un nuevo CD, el sistema permite crear y mapear zonas, pasillos y posiciones con coordenada única. (b) Cada posición tiene zona térmica asociada con rango de temperatura. (c) No se permite crear posiciones sin zona asignada. (d) Cambios de topología quedan registrados con usuario, fecha y detalle del cambio. & Alta / E1 &  \\
  RF-02.02 & Jefa de Bodega / Gerente de Operaciones & El sistema debe consolidar el stock físico en unidades de cada SKU en todos los nodos de la red logística. & Al menos un CD está configurado (RF-02.01) con inventario registrado. & El sistema expone un panel donde el usuario visualiza: (a) el stock físico totalizado por SKU, (b) el desglose por instalación, y (c) el desglose por zona operativa dentro de cada instalación. & Alta / E1 &  \\
  RF-02.03 & Jefa de Bodega & El sistema debe calcular periódicamente una propuesta de reasignación de ubicaciones (slotting) basada en la clasificación de rotación (ABC) del último período configurable, considerando además peso, volumen y compatibilidad de zona. Los productos de clasificación A deben proponerse en ubicaciones cercanas a los andenes de salida. La propuesta debe ser aprobada por el jefe de bodega antes de generar instrucciones de movimiento. & Existe inventario registrado con al menos 1 período de datos de movimiento para calcular ABC. & (a) Al ejecutar el cálculo, el sistema presenta una propuesta con: producto, ubicación actual, ubicación propuesta, motivo (clasificación ABC, peso, volumen o zona). (b) El jefe de bodega puede aprobar, rechazar o modificar cada movimiento individualmente. (c) Los movimientos aprobados generan instrucciones de reubicación asignables a dispositivos de operarios, con indicación de ubicación origen y destino. (d) Cada reubicación ejecutada queda registrada con operario, fecha/hora, y ubicaciones origen/destino. & Alta / E1 & ) \\
  RF-02.04 & Jefa de Bodega & El sistema debe generar un plan de conteo cíclico mensual que cubra una cantidad de posiciones definidas, seleccionando posiciones por criterios configurables (rotación, valor, discrepancia histórica). Las tareas de conteo deben asignarse a los dispositivos de los operarios de bodega en modo de conteo ciego, sin mostrar la cantidad teórica. Tras el ingreso de la cantidad física, el sistema compara contra el teórico y registra la diferencia con su causa en campo estructurado. & Existe inventario registrado con ubicaciones asignadas (RF-02.01). & (a) El operario recibe en su dispositivo la lista de posiciones a contar, sin cantidad teórica visible. (b) Al ingresar la cantidad física, si la diferencia valorizada es menor o igual al umbral configurado por el jefe de bodega para ese recinto, el sistema registra el ajuste automáticamente con la causa seleccionada (sobrante, falta, error de ubicación, daño, vencimiento). (c) Si la diferencia supera el umbral, el sistema marca la tarea como 'Pendiente de Autorización' y notifica al jefe de bodega. & Alta / E1 & . 14.1 \\
  RF-02.05 & Preventistas / Jefa de Bodega & El sistema debe calcular el stock disponible para venta por SKU y por centro de distribución, definido como la diferencia entre el stock físico registrado y el inventario comprometido por pedidos confirmados pendientes de preparación. En modo desconectado, la consulta debe reflejar el stock al momento de la última sincronización, indicando la antigüedad del dato. & Existe inventario registrado y pedidos confirmados. & (a) Al consultar stock, el preventista visualiza la cantidad disponible por SKU para su CD asignado. (b) La consulta refleja el stock actualizado tras cada transacción confirmada. (c) En modo offline, la consulta muestra «Dato de las HH:MM del DD/MM». (d) Al sincronizar un pedido tomado offline, si el stock disponible es insuficiente para alguna línea, el sistema marca esa línea como «sin stock disponible» y notifica al preventista. (e) Los conflictos entre pedidos simultáneos se resuelven por orden de confirmación en el servidor central (S-08, RNG-01); un pedido capturado sin señal toma su lugar al sincronizar. & Alta / E1 & ) \\
  RF-02.06 & Jefa de Bodega / Gerente de Operaciones & El sistema debe registrar como «en tránsito» todo inventario despachado desde un nodo de origen hacia un nodo destino, desde la carga confirmada hasta la recepción confirmada en destino, y generar la transacción de transferencia correspondiente en el ERP. & Existen al menos dos nodos configurados (RF-02.01) con stock registrado. & (a) Al consultar stock, el usuario visualiza: stock físico en sitio, stock comprometido, stock disponible y stock en tránsito hacia ese sitio con fecha/hora estimada de llegada. (b) El stock en tránsito NO se suma al disponible para venta hasta la recepción confirmada en destino. (c) Al confirmar la recepción en destino, el stock se transfiere automáticamente de «en tránsito» a «físico en sitio». (d) Si el nodo destino está offline al recibir el envío, el operario registra la recepción localmente y el sistema sincroniza al reconectarse. (e) Cada transferencia genera la transacción de movimiento inter-bodegas en el ERP; si el ERP no está disponible, la transacción se encola para reintento automático. (f) El stock en tránsito se muestra como disponibilidad futura, diferenciado del stock físico presente en el sitio. & Media / E1 & ) \\
  RF-02.07 & Jefa de Bodega / Gerente de Operaciones & El sistema debe permitir consultar el nivel de ocupación de cada centro de distribución, desglosado por zona operativa (seco, refrigerado, congelado), expresado como porcentaje de ubicaciones físicas utilizadas respecto de la capacidad total configurada. & CD configurado (RF-02.01) con inventario registrado. & (a) Al consultar la ocupación de un CD, el sistema muestra el porcentaje de ocupación total y por zona (b) El reporte incluye un histórico diario de al menos 12 meses para visualizar tendencias. (c) El reporte es exportable a formato de hoja de cálculo. & Alta / E1 & ). \\
  RF-02.08 & Jefa de Bodega & El sistema debe dirigir la preparación de pedidos (picking) asignando tareas de recolección al dispositivo del preparador, indicando: ubicación, producto, lote, cantidad, y secuencia optimizada de recorrido. El preparador confirma cada línea en el dispositivo. Al no encontrar un producto o encontrar cantidad insuficiente, el preparador debe registrar el faltante con su causa en un campo estructurado (sin stock, producto dañado, producto vencido, error de ubicación, otro). & Existen SKU con control de vencimiento habilitado (RF-01.03, RF-01.10) y stock con fecha de vencimiento registrada. & (a) El preparador recibe en su dispositivo la lista de líneas a recolectar, ordenadas por secuencia de recorrido óptima. (b) La confirmación de cada línea no supera 1 segundo de respuesta en el dispositivo. (c) Al registrar un faltante, el sistema exige la selección de una causa antes de continuar. (d) Al finalizar la preparación, el sistema actualiza el stock automáticamente y envía los faltantes al área comercial para gestión con el cliente. & Alta / E1 & Cap. 8, entrevista de Ximena Bravo; Cap. 15, RT-09.01; Cap. 18, criterio 15 (preparación y causa de faltantes). \\
  RF-02.09 & Jefa de Bodega & El sistema debe gestionar el almacenamiento de productos con fecha de vencimiento aplicando la regla FEFO (First Expired, First Out), priorizando en la preparación de pedidos los lotes con fecha de vencimiento más próxima. El sistema debe alertar cuando un producto en stock alcance un umbral de vida útil remanente configurable. & Existen SKU con control de vencimiento habilitado (RF-01.03, RF-01.10) y stock con fecha de vencimiento registrada en el sistema. & (a) Al generar las tareas de picking, el sistema prioriza las ubicaciones con menor vida útil remanente para cada SKU. (b) Si un producto en stock tiene una vida útil remanente inferior al umbral configurado (ej. 30 días), el sistema genera una alerta al jefe de bodega con la acción recomendada (descuento, devolución al proveedor, destrucción). (c) El sistema genera un reporte semanal de productos próximos a vencer por zona y por SKU. & Alta / E1 & Cap. 9.10; Cap. 12, inocuidad alimentaria (vencimientos y merma). \\
  RF-02.10 & Jefa de Calidad / Jefa de Bodega & El sistema debe proveer una interfaz de trazabilidad que permita buscar y detallar el ciclo de vida completo de un lote. Al consultar un lote específico, debe desglosar su estado actual, ubicación física en bodega, datos de recepción de proveedor y detalle de salida (clientes/mermas). & El lote a consultar existe en el sistema. & (a) El sistema provee un buscador de lotes por SKU, fecha o proveedor, devolviendo una lista con el estado (activo, agotado, en cuarentena). (b) Al seleccionar un lote, se despliega su ficha detallando: SKU, fecha vencimiento y cantidad física remanente por ubicación (pasillo/posición). (c) Se detalla el origen: fecha de recepción, proveedor, cantidad recibida inicial y documento. (d) Se detalla la salida de unidades: guías de despacho, cliente, fecha y cantidad entregada. (e) Se detallan excepciones: mermas, devoluciones o ajustes asociados. (f) Todo el detalle es exportable a hoja de cálculo para retiro sanitario inmediato. & Alta / E1 & . 4.9, Cap. 9.5 \\
  RF-03.01 & Preventistas & La app debe mostrar, con conectividad, el stock disponible en tiempo real para cada SKU al tomar el pedido. & El preventista tiene conectividad. & El preventista ve la disponibilidad vigente antes de ofrecer un producto. & Crítica / E1 & Cap. 9.2; Cap. 18, criterio 5. \\
  RF-03.02 & Preventistas & La app debe mostrar, sin conectividad, el stock según la última sincronización válida, con su antigüedad visible. & Sin conectividad y con sincronización previa. & El preventista ve el dato cacheado y su antigüedad. & Crítica / E1 & Cap. 9.2; Cap. 10, restricción 3 (consulta y operación sin señal). \\
  RF-03.03 & Preventistas & El sistema debe reservar el stock comprometido en el momento de la confirmación del pedido en el servidor central de la solución operacional, por orden de llegada a ese servidor (S-08, RNG-01), sin depender de la disponibilidad del ERP. & El pedido fue confirmado en el servidor central de la solución. & Si dos preventistas comprometen el mismo producto, el segundo en confirmar queda con quiebre parcial registrado. & Alta / E1 & . \\
  RF-03.04 & Preventistas & La app debe mostrar, con conectividad, el crédito disponible y la deuda vencida del cliente al tomar el pedido. & El preventista tiene conectividad. & El preventista ve crédito y deuda antes de finalizar el pedido. & Crítica / E1 & Cap. 9.2; Cap. 18, criterio 5. \\
  RF-03.05 & Preventistas & La app debe mostrar, sin conectividad, el crédito y la deuda según la última sincronización válida, con su antigüedad visible. & Sin conectividad y con sincronización previa. & El preventista ve el dato cacheado y su antigüedad. & Crítica / E1 & Cap. 9.2; Cap. 10, restricción 3 (información comercial y operación sin señal). \\
  RF-03.06 & Preventistas & La app debe mostrar el historial de compra del cliente (productos y frecuencia). & El cliente tiene al menos una compra registrada. & El preventista consulta el historial e identifica productos habituales o quiebres. & Alta / E1 & Cap. 9.2. \\
  RF-03.07 & Preventistas & La app debe permitir aplicar una promoción vigente (3x2, descuento por volumen) a una línea de pedido, calculando el descuento automáticamente. & Existe una promoción vigente para el producto y período. & El preventista aplica la promoción y el sistema calcula el descuento. & Alta / E1 & Cap. 9.2 (promociones y qué ofrecer). \\
  RF-03.08 & Preventistas & El sistema debe alertar al preventista cuando el pedido, sumado a la deuda vencida, supere el límite de crédito del cliente. & El cliente tiene un límite de crédito definido. & El preventista es alertado antes de enviar el pedido. & Crítica / E1 & Cap. 8, entrevistas de Rubén Salinas y Sandra Riffo; Cap. 9.2 (crédito y pedidos retenidos). \\
  RF-03.09 & Preventistas & El sistema debe bloquear el envío del pedido cuando este, sumado a la deuda vencida, supere el límite de crédito en más de 110 \%. & Se generó la alerta de exceso de crédito (RF-03.08). & El pedido no puede enviarse mientras el bloqueo esté activo. & Crítica / E1 & Cap. 9.2 (control de crédito). \\
  RF-03.10 & Gerente de Finanzas & El sistema debe permitir que un supervisor autorice, de forma registrada, el envío de un pedido bloqueado. & Un pedido está bloqueado (RF-03.09). & Queda registrado supervisor, fecha y hora, y el pedido puede enviarse. & Alta / E1 & ). \\
  RF-03.11 & Gerente Comercial & El sistema debe registrar y conservar en cada pedido de preventa el precio acordado por el preventista al tomarlo, conforme a las condiciones comerciales autorizadas, sin modificarlo por cambios posteriores de la lista de precios. & El preventista registra un pedido con precios y condiciones comerciales autorizados. & El pedido conserva el precio acordado y lo transmite al ERP para la facturación. & Media / E1 & . \\
  RF-03.12 & Gerente Comercial & El sistema debe generar una alerta al responsable comercial cuando el importe calculado para el despacho o la facturación difiera del precio registrado en el pedido, para corregir la discrepancia antes de facturar y conservar el precio acordado (RF-03.11). & Existe una discrepancia entre el precio registrado en el pedido y el importe calculado para el despacho o la facturación. & La diferencia queda alertada al responsable comercial y se corrige antes de facturar, conservando el precio acordado. & Media / E1 & . \\
  RF-03.13 & Preventistas & El sistema debe informar al preventista la fecha y ventana horaria de entrega que puede comprometer, considerando stock, capacidad de despacho y restricciones del cliente. & El módulo de planificación de rutas expone capacidad disponible (dependencia externa). & El preventista visualiza la ventana disponible antes de confirmar; la promesa queda registrada. & Crítica / E1 & Cap. 9.1 (promesa de fecha y franja horaria). \\
  RF-03.14 & Preventistas & El sistema debe generar sugerencias de productos usando historial de compra, frecuencia, productos habituales y promociones vigentes. & El cliente tiene historial de compra registrado. & Al seleccionar un cliente, se muestran sugerencias según criterios configurados por comercial. & Alta / E1 & ). \\
  RF-03.15 & Preventistas & La app debe generar un identificador único en el dispositivo al crear un pedido sin conectividad. & El preventista está registrando un pedido sin señal. & Cada pedido creado sin conectividad tiene un identificador único generado localmente. & Crítica / E1 & Cap. 9.2; Cap. 18, criterio 6 (pedidos sin pérdida ni duplicación). \\
  RF-03.16 & Preventistas & El sistema debe detectar y descartar automáticamente los reenvíos duplicados de un pedido con el mismo identificador único al sincronizar. & El dispositivo recupera conectividad y reintenta el envío (RF-03.16). & Ningún pedido se pierde ni se duplica al sincronizar; todo reintento queda auditado. & Crítica / E1 & Cap. 9.2; Cap. 18, criterio 6 (pedidos sin pérdida ni duplicación). \\
  RF-04.01 & Planificador de Rutas & El sistema debe asignar pedidos a la flota, secuenciando la ruta en función de la capacidad volumétrica y zonas logísticas de entrega. & Existen pedidos confirmados pendientes de preparación y despacho (RF-03.03) y la flota registrada con sus capacidades nominales de peso y volumen. & Se despliega una ruta propuesta con la carga asignada a los vehículos de la flota. & Alta / E1 & Caso, Cap 4.4, Criterio de aceptación \#7 \\
  RF-04.02 & Planificador de Rutas & El sistema debe permitir modificar manualmente la secuencia de ruta y reasignar pedidos entre vehículos & Existe una ruta generada por el módulo de planificación (RF-04.01) disponible para modificación. & La tabla de ruteo y el ETA se recalculan inmediatamente tras realizar alguna reasignación & Media / E1 & Cap. 18, criterio 7 (corrección de rutas). \\
  RF-04.03 & Planificador de Rutas & El sistema debe bloquear la adición de líneas de pedido a una ruta si la sumatoria excede los límites de peso o volumen nominales del vehículo. & El planificador intenta agregar líneas de pedido a una ruta existente (RF-04.01) y el vehículo tiene límites nominales de peso y volumen configurados. & Se despliega un error bloqueante impidiendo el cierre del camión hasta reubicar el excedente. & Alta / E1 & ). \\
  RF-04.04 & Planificador de Rutas & El sistema debe permitir la configuración de ventanas horarias específicas por cliente y ajustar la secuencia para cumplirlas. & Los clientes de la ruta tienen ventanas horarias de recepción configuradas en su ficha maestra (RF-04.07). & El planificador secuencia todas las entregas en ventanas donde el cliente está disponible & Alta / E1 & Cap. 9.1; Cap. 9.9 (ventanas de recepción y entrega). \\
  RF-04.05 & Planificador de Rutas & El sistema debe restringir la asignación de pedidos refrigerados o congelados exclusivamente a vehículos con equipo de frío. & La flota tiene registrado si cada vehículo cuenta con equipo de frío y el pedido contiene líneas con requisito de temperatura definido por tipo de almacenamiento (RF-01.10). & Se impide la adición de cadena de frío a un camión seco. & Alta / E1 & ); Cap. 9.6. \\
  RF-04.06 & Planificador de Rutas & El sistema debe crear un aviso al contacto del cliente con la hora estimada de llegada incluyendo una ventana de 30 minutos. & La ruta fue planificada y confirmada (RF-04.01) y el cliente tiene un contacto de notificación registrado en su ficha maestra. & El cliente recibe de forma exitosa el mensaje antes de la partida del camión. & Media / E1 & Cap. 15, RT-16.21 (avisos de despacho y llegada). \\
  RF-04.07 & Planificador de Rutas & El sistema debe permitir parametrizar franjas horarias de recepción en la ficha maestra de cada cliente (Ventana Horaria). & El planificador cuenta con permisos para parametrizar la ficha maestra del cliente y el cliente existe como activo en el sistema de gestión. & Ningún horario estimado de llegada es asignado en un horario fuera de la ventana horaria establecida por el cliente. & Alta / E1 & Cap. 9.1 (ventanas horarias). \\
  RF-04.08 & Gerente de Finanzas & El sistema debe calcular el costo de cada entrega tras haber sido realizada. & La entrega fue ejecutada y confirmada en el dispositivo móvil (RF-06.01) y los datos de kilómetros recorridos y tiempos de carga/descarga están disponibles (RF-14.04). & El sistema calcula el costo de cada entrega basado en kilómetros, tiempo de descarga, tamaño del pedido y reentregas. & Alta / E2 & Cap. 9.8; Cap. 18, criterio 11. \\
  RF-05.01 & Jefa de Bodega & El jefe de bodega debe poder iniciar la generación manualmente o configurar la generación automática al cierre de pedidos del turno. Las misiones deben permitir agrupación por pedido individual o por consolidación de pedidos de una misma ruta. & Existen pedidos confirmados pendientes de preparación con stock disponible (RF-02.05) y topología configurada (RF-02.01). & (a) Al generar misiones, la interfaz lista los ítems agrupados por pasillo lógico en secuencia de proximidad. (b) El sistema exige y permite confirmar cada línea digitalmente, exhibiendo: cantidad solicitada, SKU, lote asignado (FEFO) y coordenada exacta de ubicación. & Alta / E1 &  \\
  RF-05.02 & Preparadores & El sistema debe exigir la validación de cada línea recolectada escaneando el código GS1. El escaneo debe validar el SKU y capturar el Lote exacto tomado. Permite ingreso manual solo como contingencia auditable. & Existe una misión de picking asignada al dispositivo del preparador (RF-05.01). & Si el código de barras escaneado difiere del SKU o Lote instruido, el sistema bloquea la confirmación. Como mecanismo de contingencia, el sistema debe permitir el ingreso manual del código numérico, pero debe registrar obligatoriamente el evento en la bitácora como 'Ingreso Manual' para futura auditoría de calidad. & Alta / E1 & . 12 \\
  RF-05.03 & Preparadores & El sistema debe exigir la selección obligatoria de un motivo parametrizable de una lista estructurada, cada vez que se reporte una cantidad recolectada inferior a la solicitada. & El preparador está ejecutando una misión y la cantidad recolectada < cantidad solicitada. & Si la cantidad recolectada es inferior a la solicitada, el sistema despliega el campo obligatorio 'Causa de Faltante' con lista predefinida y deshabilita la acción 'Confirmar Línea' hasta su selección. & Alta / E1 & . \#15 \\
  RF-05.04 & Preparadores & El sistema debe secuenciar las líneas de la misión de preparación priorizando las zonas en el siguiente orden: seco $\rightarrow$ refrigerado $\rightarrow$ congelado, para minimizar el tiempo de exposición de los productos de frío fuera de su zona térmica. & Existe una misión de picking con productos de múltiples zonas térmicas. (RF-05.01) & Ante un pedido mixto, el sistema exige recolectar el 100\% de los SKUs de ambiente seco antes de liberar las ubicaciones correspondientes a las cámaras de frío en la interfaz del operario. & Alta / E1 &  \\
  RF-05.05 & Preparadores & El sistema debe asignar el inventario para la preparación de pedidos aplicando estrictamente la regla lógica FEFO (First Expired, First Out) según la fecha de vencimiento del lote. & Existe una misión con SKU que tienen control de vencimiento habilitado y múltiples lotes en stock. & (a) La misión dirige al operario a la coordenada del lote con fecha de vencimiento más próxima. (b) Al escanear un lote más nuevo cuando hay inventario del lote FEFO disponible, el sistema bloquea la confirmación con mensaje «Lote (código de lote) con vencimiento anterior disponible en [ubicación]». (c) Si el lote FEFO no está disponible (ej. dañado), el preparador puede seleccionar el siguiente lote registrando el motivo (lista estructurada: no encontrado, dañado, cantidad insuficiente). (d) La excepción FEFO queda registrada en bitácora de auditoría & Alta / E1 &  \\
  RF-05.06 & Jefa de Bodega & El sistema debe enviar al ERP, al cerrar una misión de preparación, el detalle consolidado por cada línea: SKU, cantidad efectivamente recolectada y lote correspondiente, vinculado al ID del pedido original. & La misión de picking ha sido completada y confirmada por el preparador & Al ejecutar 'Cerrar Misión', el sistema transmite al ERP el objeto completo de la recolección asociado al ID del pedido. El ERP procesa este conjunto de datos ajustando las cantidades finales y registrando los lotes exactos para la emisión de la Guía de Despacho, sin permitir digitación manual. & Alta / E1 & ) \\
  RF-05.07 & Jefa de Bodega & El sistema debe dirigir la carga de pedidos preparados en el camión siguiendo la secuencia inversa de la ruta asignada, garantizando que el primer cliente a visitar sea el último cargado (más accesible). El cargador debe confirmar cada bulto/pallet escaneando su SSCC. & Los pedidos están preparados y asignados a un camión con ruta definida. & (a) El dispositivo del cargador indica el siguiente bulto/pallet a subir, mostrando: SSCC, cliente destino y posición en la secuencia. (b) El cargador confirma cada bulto escaneando el SSCC; si escanea un bulto fuera de secuencia, el sistema alerta y solicita confirmación del motivo. (c) Al completar la carga, el sistema verifica que todos los bultos asignados al camión han sido escaneados. (d) Si falta algún bulto, el sistema bloquea la confirmación de carga completa y genera una alerta al jefe de bodega. & Alta / E1 &  \\
  RF-06.01 & Tripulación propia y peoneta/ Conductores externos & El sistema debe permitir marcar la recepción conforme de bultos mediante validación por línea de pedido individual o marcado masivo por guía. & El conductor tiene cargada en su dispositivo la guía de despacho electrónica del pedido (RF-05.06) y se encuentra en el local del cliente. & Los elementos marcados como entregados quedan registrados de forma exitosa en la guía. & Alta / E1 & Cap. 9.4; Cap. 15, RT-16.14 (registro y evidencia de recepción). \\
  RF-06.02 & Tripulación propia y peoneta / Conductores externos/ Canal Tradicional / Food Service / Canal Moderno) & El sistema debe permitir capturar la firma digital del receptor como evidencia de entrega, asociada a la guía de despacho electrónica. & La entrega está en curso con la guía de despacho electrónica abierta (RF-06.01) y el dispositivo cuenta con pantalla táctil funcional para la firma. & La firma queda correctamente registrada en la guía de despacho. & Crítica / E1 & Cap. 9.4; Cap. 15, RT-16.14 y RT-17.06; Cap. 16.1, decisión 15 del caso, sobre mecanismo de prueba de entrega. \\
  RF-06.03 & Tripulación propia y peoneta / Conductores externos & El sistema debe permitir reportar un rechazo total o parcial seleccionando la cantidad y el motivo desde un catálogo predefinido. & El cliente rechaza total o parcialmente la mercadería durante la entrega (RF-06.01). & Los productos rechazados quedan registrados de forma correcta en la solicitud de rechazo. & Media / E1 & Cap. 9.4; Cap. 16.1, decisión 12 (devoluciones y efectos tributarios). \\
  RF-06.04 & Tripulación propia y peoneta / Conductores externos & El sistema debe permitir reportar un local cerrado, indicando la dirección del cliente. & El conductor llega al local del cliente y este se encuentra cerrado o sin persona disponible para recibir la entrega. & El envío se registra como fallido y se agenda para ser entregado en la ventana horaria siguiente mas cercana del cliente. & Media / E1 & Cap. 16.1, decisión 3 (local cerrado). \\
  RF-06.05 & Tripulación propia y peoneta / Conductores externos & El sistema debe permitir digitar el monto exacto recaudado en pagos realizados en efectivo & El cliente paga en efectivo la entrega y el conductor se encuentra registrando la entrega correspondiente en la aplicación de reparto (RF-06.01). & El monto recibido en efectivo por el conductor queda registrado en el sistema. & Crítica / E1 & Cap. 4.7; Cap. 9.7. \\
  RF-06.07 & Tripulación propia y peoneta / Conductores externos / Jefa de Bodega / Gerente de Operaciones & El sistema debe sincronizarse automaticamente, al detectar la reconexión de señal o el regreso al centro de distribución, sincronizar automaticamente todos los registros generados durante el turno, recuperando la totalidad de los despachos registrados, consolidandolos en un reporte de cierre. & El sistema central está disponible para recibir los datos. & Al llegar al CD o recuperar señal, la app sube la data y el sistema genera automáticamente el reporte de rendición. & Crítica / E1 &  \\
  RF-06.08 & Tripulación propia y peoneta / Conductores externos / Empresas Transp. & El sistema debe permitir al conductor de una empresa externa iniciar su jornada usando las credenciales temporales por OTP (código de un solo uso) emitido por su jefe de flota; la sesión común posterior la provee RF-15.02. sin necesidad de correo corporativo. & El conductor tiene instalada la aplicación de reparto en su dispositivo, su jefe de flota le generó un código OTP vigente. Sin señal, rige la identificación por PIN o QR de RF-14.05 y la validación se completa al sincronizar. & Un conductor nuevo de una empresa externa se autentica en 2 minutos con un código que le da su jefe de flota, y puede iniciar su ruta. & Alta / E1 & Caso, Restricción \#6, RT-12.11, Cap. 16.1 (Decisión \#5) \\
  RF-06.09 & Tripulación propia y peoneta / Conductores externos & Al registrar una entrega, el sistema debe permitir al conductor capturar la cantidad de canastillos y pallets devueltos por el cliente, identificando el tipo y estado. & El conductor está en el momento de la entrega, El dispositivo del conductor puede capturar la información. & El conductor registra que el cliente devolvió 5 canastillos; el sistema descuenta del saldo del cliente y actualiza el inventario de envases. & Alta / E1 & . 9.7, Criterio Aceptación \#12 \\
  RF-06.10 & Tripulación propia y peoneta / Conductores externos & El sistema debe actualizar el saldo del cliente y del parque retornable en línea o en modo diferido según conectividad, asociando el registro a la entrega o a la guía correspondiente & Se ha integrado la devolución en el sistema (RF-06.09). & El saldo del cliente indica que se devolvieron 5 canastillos y el sistema lo actualiza en el inventario de envases & Alta / E1 & . 9.7, Criterio Aceptación \#12 \\
  RF-06.11 & Tripulación propia y peoneta / Conductores externos & El sistema debe ofrecer al cliente la opción de confirmar la recepción mediante un código QR único generado por la app del conductor, cuando el cliente disponga de teléfono inteligente. & El sistema ha generado el código QR asociado a esa guía/entrega, La guía de despacho electrónica ya está emitida & Se resuelve alternativamente la Decisión \#15; los clientes con smartphone pueden validar su recepción de forma autónoma. & Media / E1 & ) \\
  RF-06.12 & Canal Tradicional / Food Service / Canal Moderno / Tripulación propia y peoneta / Conductores externos & El sistema debe permitir confirmar la recepción del cliente aun si este no dispone de teléfono o conectividad, mediante captura de firma en pantalla o el registro fotografico de la recepcion, dejando evidencia vinculada a la guia de despacho electrónica en su acuse de recibo. & La guia de despacho esta emitida, el cliente puede firmar en pantalla o autorizar la fotografia & Los clientes sin conectividad y sin smartphone pueden validar su recepción de forma autónoma. & Media / E1 & ) \\
  RF-06.13 & Tripulación propia y peoneta / Canal Tradicional & La app de reparto debe permitir imprimir en terreno un comprobante térmico de entrega para el cliente usando una impresora portátil Bluetooth en aquellos casos donde el cliente requiera respaldo físico. & El cliente solicita un comprobante físico, El conductor tiene una impresora térmica Bluetooth conectada y con papel disponible, La impresora está emparejada con el dispositivo y la app tiene permiso para imprimir. & El cliente del canal tradicional recibe un comprobante físico pequeño, igual que hoy con la guía, pero generado digitalmente. & Media / E1 & Caso, RT-17.06 (impresora térmica portátil) \\
  RF-07.01 & Gerente de Finanzas / Tripulación propia y peoneta / Conductores externos & El sistema debe generar la planilla digital de rendición individual al retorno del camión, contrastando los valores declarados por el conductor en la aplicación móvil contra las entregas al contado realizadas y facturas cobradas. & El camión retornó a base o el dispositivo del conductor recuperó conectividad con sus registros de turno sincronizados (RF-06.07), incluyendo las entregas al contado y facturas cobradas del turno. & Despliegue automático de la cuadratura en caja mostrando monto esperado vs recaudado y la diferencia resultante a justificar. & Alta / E1 &  \\
  RF-07.02 & Gerente de Finanzas / Tripulación propia y peoneta / Conductores externos & El sistema debe obligar al registro estructurado y categorizado de causales para cualquier diferencia detectada en la rendición de caja (descuentos comerciales autorizados, diferencias de cambio, siniestro en ruta), bloqueando el cierre de la liquidación si existe saldo sin justificar. & Existe una diferencia detectada en la rendición de caja individual del camión sin justificar, generada por la cuadratura de RF-07.01. & Cero pesos de descuadre sin motivo formal tipificado, firma digital y asignación de responsabilidad en bitácora. & Alta / E1 &  \\
  RF-07.03 & Gerente de Finanzas / Canal Tradicional / Canal Moderno & El sistema debe proveer una consola de gestión para la cartera de crédito a 30 días sincronizada con el ERP, asignando carteras de clientes a ejecutivos de cobranza y registrando bitácoras de gestión, compromisos de pago y fechas de recontacto. & El ERP ha sincronizado las cuentas por cobrar de la cartera a 30 días y existen ejecutivos de cobranza configurados en el sistema. & Listado diario de morosidad al día con registro auditable de promesas de pago y tareas automáticas de seguimiento. & Alta / E1 &  \\
  RF-07.04 & Preventistas / Gerente de Finanzas / Gerente Comercial & El sistema debe notificar en la aplicación móvil de preventa el estado de bloqueo o mora de un cliente (>30 días o cupo excedido), bloqueando preventivamente el ingreso de pedidos a crédito y permitiendo únicamente venta al contado o cobro de deuda. & El cliente presenta estado de mora (>30 días) o cupo de crédito excedido informado por el ERP y el preventista selecciona al cliente en la aplicación de preventa. & Alerta en terminal de preventa en menos de 2 s al seleccionar al cliente, impidiendo el ingreso de órdenes a crédito no autorizadas. & Alta / E1 &  \\
  RF-07.05 & Tripulación propia y peoneta / Conductores externos / Canal Tradicional / Canal Moderno / Gerente de Finanzas & El sistema debe integrarse vía Bluetooth con terminales de pago portátil (POS) en la aplicación móvil de reparto para cobro con tarjeta en ruta, capturando el código de autorización bancario y conciliando el voucher automáticamente en la liquidación diaria. & El conductor dispone de un terminal POS portátil emparejado por Bluetooth con el dispositivo de reparto (RT-17.06) y el cliente paga con tarjeta en la entrega. & Captura automática del código de autorización POS y conciliación contra la entrega en la rendición diaria de caja sin digitación manual. & Alta / E1 & Caso, RT-17.06, Cap. 9.7 \\
  RF-07.06 & Canal Tradicional / Canal Moderno / Gerente de Finanzas & El sistema debe disponer en el portal web de autoatención el estado de cuenta y la cartola de antigüedad de deuda (tramos corriente, 30, 60, 90+ días) descargable en formatos PDF y Excel para clientes autenticados. & El cliente tiene una sesión activa en el portal (RF-12.26) y el ERP ha sincronizado su saldo vivo y cartola de antigüedad de deuda. & Acceso seguro del cliente autenticado a su saldo vivo y detalle de facturas pendientes con opción de descarga de cartola. & Alta / E2 &  \\
  RF-07.07 & Gerente de Finanzas / Canal Moderno & El sistema debe proveer una bandeja de validación para pagos reportados en el portal web de clientes, permitiendo al cajero verificar el comprobante bancario y autorizar su traspaso estructurado al ERP. & Existe al menos un pago reportado por un cliente en el portal web en estado «En Revisión» y el cajero liquidador tiene permisos de validación. & Pago web queda en estado «En Revisión» hasta que el cajero lo aprueba, momento en que se envía la rebaja contable al ERP. & Alta / E2 &  \\
  RF-07.08 & Jefe de TI / Gerente de Finanzas & El sistema debe transferir de forma estructurada, asíncrona e idempotente los comprobantes de recaudación y cierres comerciales aprobados hacia el ERP para actualizar automáticamente las cuentas por cobrar. & Existen comprobantes de recaudación y cierres comerciales aprobados pendientes de transferencia (RF-07.02, RF-07.03, RF-07.08) y la interfaz de integración con el ERP está configurada. & Abonos y cancelaciones contables reflejados en el ERP el mismo día sin intervención manual de digitación (Restricción R4). & Alta / E1 &  \\
  RF-07.09 & Tripulación propia y peoneta / Conductores externos / Canal Tradicional/ Gerente de Finanzas & El sistema debe permitir el registro de anticipos o pagos parciales en entregas al contado desde la aplicación móvil de reparto, emitiendo un comprobante digital de abono y traspasando el saldo insoluto a la cuenta corriente del cliente en el ERP. & El conductor está registrando una entrega al contado (RF-06.05) y el cliente realiza un anticipo o pago parcial de la deuda o la entrega en curso. & Registro de abono parcial en app móvil, comprobante digital emitido al cliente y saldo pendiente registrado en cuenta corriente. & Alta / E1 &  \\
  RF-07.10 & Preventistas / Gerente de Finanzas / Canal Tradicional & El sistema debe permitir al preventista registrar la cobranza en efectivo de facturas vencidas en el local del cliente desde su terminal móvil, emitiendo un recibo provisional digital y forzando su rendición diaria en caja. & El preventista se encuentra en el local del cliente con su terminal móvil de preventa y el cliente tiene facturas vencidas cobrables en efectivo. & Emisión de recibo digital de cobro en preventa y consolidación automática en la bolsa de dinero a rendir por el vendedor al volver a base. & Alta / E1 & . 8 \\
  RF-08.01 & Tripulación propia y peoneta / Conductores externos & El sistema debe permitir al conductor registrar la devolución de productos en el momento de la entrega, capturando SKU, cantidad y motivo (daño / vencimiento / producto no pedido / error), operando en modo offline. & El conductor está registrando la entrega en el local del cliente (RF-06.01) y el cliente devuelve productos. & La devolución queda registrada en el dispositivo al momento de la entrega, sin papeleo, y se sincroniza con el sistema central al reconectarse. & Crítica / E1 & . 4.8, RT-17.01 \\
  RF-08.02 & Jefa de Bodega & Al recepcionar una devolución en el andén del CD, el sistema debe permitir seleccionar el destino final del producto (reingreso a stock / venta con descuento / devolución al proveedor / destrucción) y enviar la transacción al ERP mediante su interfaz de integración. & Existe una devolución registrada en ruta (RF-08.01) y recepcionada en el andén del CD por el jefe de bodega. & El destino queda registrado y la transacción llega al ERP; toda decisión queda trazada con usuario, fecha y destino asignado. & Crítica / E1 & ) \\
  RF-08.03 & Tripulación propia y peoneta / Conductores externos & El sistema debe mantener una cuenta corriente de envases retornables por cliente, actualizando el saldo en cada entrega y devolución registrada por el conductor. & Existe el parque de envases retornables registrado (canastillos y pallets) y el cliente tiene una cuenta corriente de envases creada en el sistema. & El saldo de envases por cliente queda actualizado tras cada movimiento y es consultable desde el dispositivo del conductor antes de cada entrega. & Alta / E1 & . 16.1 (Decisión \#10), Cap 9.7 \\
  RF-08.04 & Preventistas & El sistema debe generar una alerta al preventista cuando el saldo de envases retornables de un cliente supere el umbral configurado por el administrador. & El cliente tiene saldo de envases retornables registrado (RF-08.03) y el administrador ha configurado el umbral máximo de saldo. & El preventista recibe la notificación en $\le$ 24 horas de superado el umbral. & Alta / E1 &  \\
  RF-08.05 & Preparadores / Preventistas & El sistema debe permitir registrar mermas por vencimiento detectadas en conteo cíclico, preparación de pedidos o góndola del cliente, asociando cada registro al SKU y al lote correspondiente. & Existe stock con fecha de vencimiento registrada (RF-01.03) y se detecta una merma por vencimiento en conteo cíclico (RF-02.04), preparación de pedidos (RF-05.03) o góndola del cliente. & La merma queda registrada por SKU y lote, con su origen (conteo, preparación o góndola), y descuenta el stock disponible. & Alta / E1 & ) \\
  RF-08.06 & Gerente de Operaciones / Gerente de Finanzas & El sistema debe reportar la pérdida de envases retornables desglosada por cliente, tipo de envase y zona, con actualización diaria. & Existen registros de movimientos y pérdidas de envases retornables por cliente, tipo de envase y zona (RF-08.03, RF-06.09, RF-06.10). & El reporte identifica los clientes y zonas con mayor pérdida de envases, permitiendo priorizar acciones de recuperación. & Alta / E1 & ) \\
  RF-08.07 & Gerente de Finanzas & El sistema debe incluir el costo de devoluciones y mermas por cliente como componente del cálculo del costo de servir (Épica 11). & El módulo de costo de servir (RF-11.02) está operativo y existen devoluciones y mermas registradas por cliente en el período (RF-08.02, RF-08.05). & El costo de servir de un cliente con alta tasa de devolución refleja ese costo adicional. & Alta / E2 &  \\
  RF-09.01 & Jefa de Calidad & El sistema debe permitir la trazabilidad 'hacia adelante' (forward tracing): dado un lote específico, debe listar todos los clientes que recibieron productos de ese lote, con fecha de entrega, guía y cantidad. & El lote a consultar existe en el sistema con su ciclo de vida registrado, incluyendo recepción y salidas a clientes (RF-02.10). & En un simulacro de retiro sanitario, se obtiene en menos de 2 horas la lista de clientes afectados, con trazabilidad desde el proveedor hasta los clientes y archivo exportable (RNF-09.01). & Crítica / E1 & . 9.6, Criterio Aceptación \#1 \\
  RF-09.02 & Jefa de Calidad & El sistema debe permitir la trazabilidad 'hacia atrás' (backward tracing): dado un producto en el local del cliente (SKU + lote), debe identificar el origen (proveedor, fecha de recepción y orden de compra) & El producto (SKU + lote) consultado existe en el sistema y su recepción de origen quedó registrada con proveedor, fecha y OC (RF-02.10). & El equipo de calidad puede rastrear el origen de cualquier unidad devuelta, realizando la trazabilidad desde los clientes afectados hasta los productores & Crítica / E1 & . 9.6 \\
  RF-09.03 & Jefa de Calidad & El sistema debe registrar de forma automática y continua la temperatura de las cámaras de refrigerado y congelado, con lecturas cada 5 minutos como mínimo & Los sensores de temperatura de las cámaras de refrigerado y congelado están instalados, conectados al sistema y enviando lecturas. & El sistema registrara las temperaturas. Los registros son auditable y se puede exportar para la autoridad sanitaria. & Crítica / E1 & . 7.3, Criterio Aceptación \#3 \\
  RF-09.04 & Jefa de Calidad & El sistema debe registrar de forma automática y continua la temperatura en los camiones con equipo de frío (18 propios y 10 de transportistas, todos con sensor instalado antes de la ola de reparto; S-28), durante todo el trayecto desde la carga hasta la entrega. & El camión con equipo de frío está equipado con sensor de temperatura conectado al sistema y el viaje está en curso desde la carga hasta la entrega. & El sistema registrara las temperaturas de los camiones. El cliente o Puelche pueden consultar el registro histórico de temperatura de su envío. & Crítica / E1 & . 4.9, Criterio Aceptación \#3 \\
  RF-09.05 & Jefa de Calidad & El sistema debe generar una alerta en tiempo real (notificación en cabina y en plataforma de monitoreo) cuando se detecte una excursión de temperatura que supere la desviación continua de 2\,°C por 15 minutos para el tipo de producto. & El sistema recibe lecturas continuas de temperatura y el umbral de excursión térmica (desviación de 2\,°C por 15 minutos) está parametrizado para el tipo de producto. & El sistema al registrar una temperatura fuera del rango aceptable, generara una alerta tanto al conductor como al equipo de calidad & Crítica / E1 & ), Cap. 16.1 (Decisión \#4) \\
  RF-09.06 & Tripulación propia y peoneta / Conductores externos & Ante una excursión térmica crítica y sostenida, el sistema debe bloquear preventivamente el despacho del lote afectado y notificar a Calidad; ante una excursión menor y transitoria, debe emitir una alerta preventiva en cabina. El conductor recibe la alerta y registra sus observaciones; Calidad decide liberar, bloquear o rechazar (LafroX, 2026c, sección «Arbitraje de tensiones operacionales y comerciales», tensión 2; Anexo 2.2, S-04). & Existe una alerta de excursión térmica activa (RF-09.05) y hay un camión con carga de productos con requisito de frío pendiente de despacho. & El lote queda retenido automáticamente ante excursión crítica; el conductor ve la alerta y registra observaciones; solo Calidad libera, bloquea o rechaza, y su decisión queda registrada. & Media / E1 & ), Cap. 16.1 (Decisión \#4) \\
  RF-09.07 & Tripulación propia y peoneta / Conductores externos & El sistema debe permitir al conductor registrar el rechazo de un producto entregado por 'sospecha de ruptura de cadena de frío', asociando el evento al sensor de temperatura del camión en ese tramo. & El cliente rechaza el producto por sospecha de ruptura de cadena de frío durante la entrega (RF-06.03) y el camión cuenta con registro del sensor en el tramo. & Si se rechaza el pedido, el conductor registra el rechazo ingresando evidencia & Crítica / E1 & . 8 \\
  RF-09.08 & Preparadores & El sistema debe registrar, en el momento de la preparación (picking), el lote específico que se le asigna a cada línea de pedido, para los productos con control sanitario. & Existe una misión de picking en ejecución con líneas de productos con control sanitario. & El picking digital registra qué lote concreto sale en cada canastillo/pallet hacia el cliente. & Crítica / E1 & . 4.9 \\
  RF-09.09 & Jefa de Calidad & El sistema debe generar la evidencia de cumplimiento de cadena de frío exigible ante la autoridad sanitaria y ante el cliente, para un envío determinado. & El envío fue entregado y existen los registros de temperatura del trayecto (RF-09.03, RF-09.04) junto con los datos de trazabilidad del envío (RF-09.01). & La evidencia de cadena de frío de un envío se obtiene sin reconstrucción manual. & Crítica / E1 & Cap. 18 criterio N°3 \\
  RF-11.01 & Gerente de Finanzas / Gerente Comercial / Jefa de Calidad & El sistema debe programar y distribuir vía correo electrónico reportes ejecutivos consolidados (en formatos PDF y XLSX) con periodicidad semanal y mensual, desglosando la evolución de OTIF, Fill Rate, Merma por vencimiento y Costo de Servir por zona y canal de comercialización. & Existen datos calculados de OTIF, Fill Rate, merma por vencimiento y costo de servir del período (RF-11.02, RF-11.03, RF-11.04) y los destinos de correo de los gerentes están configurados. & Reportes emitidos automáticamente a las casillas de los gerentes el día 2 de cada mes y cada lunes a las 06:00 AM sin intervención manual. & Media / E2 &  \\
  RF-11.02 & Gerente de Finanzas / Gerente Comercial & El sistema debe calcular el Costo de Servir real por entrega y por cliente, imputando los costos directos de transporte (combustible devengado vía GPS, peajes, tiempo de conducción y tiempo de descarga) y los costos indirectos prorrateados (bodega y administración). & Existen entregas ejecutadas con sus costos directos asociados (telemetría RF-14.04, kilómetros y tiempos de descarga) y los costos indirectos de bodega y administración están definidos. & Ficha de rentabilidad por cliente actualizada tras cada liquidación, desglosando costos logísticos directos e indirectos y margen neto resultante. & Media / E2 &  \\
  RF-11.03 & Gerente Comercial / Jefa de Calidad / Tripulación propia y peoneta / Conductores externos & El sistema debe calcular diariamente el indicador OTIF por cliente, ruta, zona y total consolidado, definiendo «On-Time» como el cumplimiento de fecha y ventana horaria pactada e «In-Full» como la entrega del 100\% de ítems y unidades pedidas, obligando al registro de causal tipificada ante desvíos. & Existen pedidos confirmados con ventana horaria pactada (RF-04.04, RF-04.07) y entregas registradas con su cumplimiento y POD (RF-06.01). & Todas las gerencias visualizan el mismo valor de OTIF diario; ante entregas parciales o tardías se audita el motivo en el sistema. & Alta / E1 &  \\
  RF-11.04 & Gerente Comercial / Tripulación propia y peoneta / Conductores externos & El sistema debe calcular automáticamente el Fill Rate (unidades entregadas vs. pedidas) y el índice de Perfect Order contrastando la orden de venta original contra los datos de entrega física confirmados en el dispositivo móvil de reparto, sin requerir interacción tecnológica del cliente tradicional. & Existen órdenes de venta originales y sus entregas físicas confirmadas en el dispositivo móvil de reparto (RF-06.01). & Indicador mensual de Perfect Order desglosado por causa raíz (daño, faltante, atraso, documento) contrastado contra la orden de venta original. & Alta / E1 &  \\
  RF-11.05 & Planificador de Rutas / Gerente Comercial & El sistema debe calcular diariamente la utilización de capacidad volumétrica (m\textasciicircum{}3) y de peso (kg) por viaje despachado, alertando aquellos camiones que salgan con una ocupación inferior al umbral configurable (ej. < 65\%). & Existen rutas despachadas con sus datos de carga en m³ y kg (RF-04.01, RF-04.03) y el umbral configurable de ocupación está definido. & Generación de bitácora de viajes subutilizados en la ventana de 05:30 a 07:00 para auditoría y balanceo de carga de transporte. & Alta / E1 &  \\
  RF-11.06 & Gerente Comercial / Gerente de Finanzas / Jefa de Calidad & El sistema debe proveer un Tablero de Control Operacional en tiempo real que consolide el avance de rutas de despacho, porcentaje de entregas cumplidas, devoluciones en tránsito, alertas activas de cadena de frío y monto recaudado en efectivo. & Existen datos en tiempo real de rutas en ejecución, entregas, devoluciones, alertas de cadena de frío y recaudación (RF-04.01, RF-06.05, RF-06.07, RF-09.05). & Despliegue en vivo en pantallas gerenciales con actualización continua de la distribución y alertas críticas de ruta. & Alta / E1 & Caso, RT-05.29, Cap. 9.1, 9.6 \\
  RF-11.07 & Gerente de Finanzas / Tripulación propia y peoneta / Conductores externos & El sistema debe procesar automáticamente la liquidación y cierre comercial individual por camión al retornar a base, consolidando documentos tributarios entregados, devoluciones físicas recibidas, cobranzas en efectivo y balance de envases retornables. & El camión retornó a base con los registros del turno sincronizados (RF-06.07): documentos tributarios entregados, devoluciones, cobranzas en efectivo y balance de envases retornables. & Planilla de cuadratura emitida inmediatamente al arribo del camión para tesorería y bodega, sin rezagos para el día siguiente. & Alta / E1 & Caso, RT-05.29, CA-10 \\
  RF-11.08 & Gerente Comercial / Preventistas / Gerente de Finanzas & El sistema debe segmentar mensualmente la cartera de clientes en matrices de rentabilidad neta (margen comercial vs. costo de servir real), sugiriendo ajustes de frecuencia de visita o umbrales mínimos de compra para clientes no rentables. & Existen datos de margen comercial y costo de servir por cliente (RF-11.02) de al menos un período mensual cerrado. & Clasificación automática de clientes en rentables, no rentables y estratégicos, con simulación de impacto de cambio de frecuencias. & Media / E2 &  \\
  RF-12.01 & Canal Moderno & El sistema debe recibir el mensaje de pedido electrónico estructurado (EDI/EDIFACT/XML, según el estándar que exija cada cadena registrada del canal moderno). & La cadena está registrada con su mapeo de mensajería configurado. & El mensaje recibido queda disponible para su validación (RF-12.02) en un plazo no superior a 15 segundos, identificando la cadena de origen. & Alta / E2 & RT-05.23 \\
  RF-12.02 & Canal Moderno & El sistema debe validar automáticamente el pedido EDI recibido, verificando estructura del mensaje, existencia del cliente, existencia de los productos, cantidades, precios y demás datos obligatorios de la cadena. & El mensaje fue recibido (RF-12.01). & El pedido queda marcado como válido o inválido, con el detalle de cada validación fallida. & Alta / E2 & . 17.1 \\
  RF-12.03 & Canal Moderno & El sistema debe resolver el código de producto informado por la cadena contra el maestro de productos vigente de Puelche, mediante una tabla de equivalencias por cadena. & El pedido pasó la validación de estructura (RF-12.02). & Cada línea queda asociada a un SKU válido del maestro interno, o marcada como discrepancia no resuelta. & Alta / E2 &  \\
  RF-12.04 & Canal Moderno & El sistema debe registrar automáticamente, sin digitación manual, el pedido EDI válido (RF-12.02, RF-12.03) en el sistema de gestión empresarial. & El pedido fue validado y sus productos resueltos contra el maestro. & El pedido aparece como pedido único en el sistema de gestión dentro de 2 minutos desde su validación. & Alta / E2 &  \\
  RF-12.05 & Planificador de Rutas & El sistema debe poner a disposición del proceso de planificación de rutas todo pedido registrado proveniente de un canal electrónico. & El pedido fue registrado en el sistema de gestión (RF-12.04). & El pedido aparece en el listado de pedidos confirmados del módulo de planificación, sin digitación manual. & Alta / E2 &  \\
  RF-12.06 & Gerente Comercial & El sistema debe permitir gestionar, mediante una bandeja de excepciones, los pedidos que no pasaron la validación (RF-12.02) o la resolución de maestro (RF-12.03), permitiendo corregir, rechazar o reprocesar cada uno. & Existe al menos un pedido en estado de excepción. & Los pedidos con error aparecen en la bandeja con causa, estado, fecha y responsable. & Alta / E2 &  \\
  RF-12.07 & Canal Moderno & El sistema debe enviar a la cadena de origen una confirmación o rechazo del pedido, indicando la causa cuando no pueda ser aceptado completamente. & El pedido fue validado (RF-12.02) o resuelto en la bandeja de excepciones (RF-12.06). & La cadena recibe automáticamente el mensaje correspondiente. & Alta / E2 &  \\
  RF-12.08 & Gerente Comercial / Planificador de Rutas & El sistema debe mantener el historial de estados de cada pedido electrónico (recepción, validación, confirmación, preparación, despacho, ASN y prueba de entrega). & El pedido fue recibido (RF-12.01). & Se puede consultar el estado actual y el historial de eventos con fecha, hora y responsable/sistema & Alta / E2 &  \\
  RF-12.09 & Canal Moderno & El sistema debe generar y enviar al cliente del canal moderno un ASN cuando su pedido queda cargado en el camión, referenciando la guía de despacho electrónica ya emitida, con productos, cantidades y hora estimada de llegada. & El pedido fue cargado y la guía de despacho electrónica fue emitida. & El cliente recibe el ASN dentro de los 2 minutos posteriores a la carga. & Alta / E2 & entrevista Rubén Salinas \\
  RF-12.10 & Planificador de Rutas & El sistema debe planificar la llegada del camión al cliente del canal moderno dentro de la ventana horaria de 30 minutos comprometida. & Existe un módulo de ruteo capaz de programar por ventana horaria. & El camión llega dentro de la ventana comprometida en el 80 \% de los casos. & Alta / E2 &  \\
  RF-12.11 & Tripulación propia y peoneta / Conductores externos & El sistema debe registrar el cumplimiento o incumplimiento de la ventana de 30 minutos para cada entrega del canal moderno, distinguiéndola de la ventana crítica de despacho 05:30-07:00. & La entrega fue ejecutada o el plazo comprometido venció. & Toda desviación queda registrada con hora real de llegada. & Alta / E2 & RT-10.05 \\
  RF-12.12 & Tripulación propia y peoneta / Conductores externos & El sistema debe capturar evidencia de entrega en el local del cliente: firma en pantalla del titular cuando pueda firmar, o fotografía de la mercadería más nombre de quien recibe cuando no sea el titular o no pueda firmar. & La entrega fue realizada. & Queda un único registro de evidencia por entrega, con el mecanismo utilizado identificado. & Alta / E2 & Cap. 16.1, Decisión \#15 \\
  RF-12.13 & Canal Moderno & El sistema debe enviar al cliente del canal moderno un acuse de recibo digital de la entrega, articulado con la guía de despacho electrónica y su acuse de recibo tributario. & La evidencia fue capturada (RF-12.12) y la guía electrónica fue emitida. & El cliente recibe el comprobante dentro de los 30 minutos posteriores a la descarga. & Alta / E2 &  \\
  RF-12.14 & Canal Tradicional / Food Service / Canal Moderno & El sistema debe publicar un portal web público, sin autenticación, con el catálogo vigente de productos (sin precios ni stock), contacto comercial e información institucional. & Un visitante accede al portal público para consultar el catálogo de productos sin autenticación (). & Cualquier visitante ve el catálogo sin poder generar pedidos ni ver precios. & Media / E2 & Caso, RT-16.30 \\
  RF-12.15 & Canal Moderno & El portal autenticado de clientes debe permitir generar un pedido en autoservicio, validando el mismo crédito y stock que consulta el preventista. & El cliente tiene una sesión activa (RF-12.26). El cliente existe como cliente activo del canal moderno en el sistema de gestión. & El cliente genera un pedido válido sin intervención de un preventista. & Media / E2 & Caso, RT-16.30 \\
  RF-12.16 & Canal Moderno & El portal de clientes debe permitir consultar el estado de sus entregas y la ventana de entrega comprometida. & El cliente tiene una sesión activa (RF-12.26) y tiene al menos un pedido registrado a su RUT. & El cliente visualiza el estado y compromiso de cada pedido. & Media / E2 & Caso, RT-16.30 / Cap. 9.1 \\
  RF-12.17 & Canal Moderno & El portal de clientes debe permitir consultar y descargar sus documentos tributarios (factura, guía electrónica, nota de crédito). & El cliente tiene una sesión activa (RF-12.26) y el documento fue emitido a su RUT por el sistema de gestión. & El cliente descarga los documentos correspondientes a sus operaciones. & Media / E2 & Caso, RT-16.30 / RT-16.14 \\
  RF-12.18 & Canal Moderno & El portal de clientes debe permitir consultar el saldo y estado de cuenta del cliente autenticado, con datos del sistema de gestión. & El cliente tiene una sesión activa (RF-12.26). & El cliente consulta su saldo según la información del sistema de gestión. & Media / E2 & Caso, RT-16.30 \\
  RF-12.19 & Empresas Transp. & El portal de transportistas debe permitir consultar las rutas asignadas a la empresa para el día siguiente. & La empresa está registrada y tiene rutas asignadas. & La empresa visualiza el día anterior las rutas asignadas. & Media / E2 & Caso, RT-16.30 \\
  RF-12.20 & Empresas Transp. & El portal de transportistas debe permitir descargar la documentación de despacho de cada ruta asignada. & La ruta fue asignada y la documentación fue generada. & La empresa descarga la documentación correspondiente. & Media / E2 & Caso, RT-16.30 \\
  RF-12.21 & Empresas Transp. & El portal debe permitir confirmar, antes de la ventana de despacho, qué conductor y vehículo ejecutarán cada ruta asignada del día. & La ruta fue asignada a la empresa transportista. & Queda registrado el conductor/vehículo real; si no se confirma, el sistema mantiene la asignación planificada como supuesto no verificado. & Media / E2 &  \\
  RF-12.22 & Proveedores & El portal de proveedores debe permitir consultar, en modo de solo lectura, el estado de sus órdenes de compra. & El proveedor está registrado y tiene al menos una orden vigente. & El proveedor consulta el estado actualizado de sus órdenes. & Media / E2 & Caso, RT-16.30 \\
  RF-12.23 & Proveedores & El portal de proveedores debe permitir consultar, en modo de solo lectura, las recepciones registradas para sus órdenes. & Existe al menos una recepción registrada. & El proveedor consulta las recepciones actualizadas. & Media / E2 & Caso, RT-16.30 \\
  RF-12.24 & Proveedores & El portal de proveedores debe permitir consultar, en modo de solo lectura, sus devoluciones y notas de crédito. & Existe al menos una devolución o nota de crédito registrada. & El proveedor consulta las devoluciones/notas de crédito actualizadas. & Media / E2 & Caso, RT-16.30 \\
  RF-12.25 & Canal Moderno & El portal debe permitir al cliente del canal moderno registrarse con su RUT de empresa, correo electrónico corporativo y una contraseña, siempre que el RUT esté previamente habilitado en el sistema de gestión como cliente activo de Puelche. & El cliente existe como cliente activo en el sistema de gestión (2017) y no tiene cuenta creada en el portal. La forma de identificar al cliente en ese sistema se consulta en el Anexo 3.H. & El cliente queda registrado y recibe un correo de confirmación de cuenta. Si el RUT no está habilitado, el sistema informa que debe contactar a su ejecutivo comercial. & Media / E2 & ). \\
  RF-12.26 & Canal Moderno & El portal debe permitir al cliente autenticarse con su correo y contraseña registrados. & El cliente tiene una cuenta creada y confirmada (RF-12.25). & El cliente accede a su área privada. Tras tres intentos fallidos consecutivos, la cuenta queda bloqueada temporalmente. & Media / E2 & ). \\
  RF-12.27 & Canal Moderno & El portal debe permitir al cliente solicitar un enlace de restablecimiento de contraseña a su correo registrado cuando no recuerde su clave. & El cliente tiene una cuenta creada (RF-12.25) y acceso al correo registrado. & El cliente recibe un enlace válido por tiempo limitado que le permite establecer una nueva contraseña. El enlace expira tras su uso o transcurridos 15 minutos. & Media / E2 & ). \\
  RF-12.28 & Canal Moderno & El portal debe permitir al cliente cerrar su sesión de forma explícita, invalidando el token de sesión activo. & El cliente tiene una sesión iniciada (RF-12.26). & La sesión queda invalidada y el cliente es redirigido a la pantalla de inicio de sesión. & Media / E2 & ). \\
  RF-12.29 & Canal Moderno & El portal debe cerrar automáticamente la sesión del cliente tras un período de inactividad definido, sin requerir acción explícita de cierre. & El cliente tiene una sesión iniciada (RF-12.26) y no ha interactuado con el portal durante el período configurado. & La sesión expira y el cliente debe iniciar sesión nuevamente para acceder a su área privada. & Media / E2 & ). \\
  RF-12.30 & Canal Moderno & El portal debe permitir al cliente pagar los pedidos generados en autoservicio (RF-12.15) mediante la plataforma de pagos Transbank. & El cliente tiene un pedido confirmado y sesión activa (RF-12.26). & El pago es procesado por Transbank y el portal registra la confirmación o rechazo de la transacción. & Media / E2 & ). \\
  RF-14.01 & Gerente de Operaciones & El sistema debe integrarse con la plataforma de telemetría de terceros existente que cubre los camiones propios para consumir datos de posición GPS, velocidad y kilometraje. & Existe un contrato vigente con el proveedor de telemetría para los 42 camiones propios, El sistema central tiene conectividad hacia la plataforma del proveedor (Internet o red privada), se ha obtenido la credencial tecnica y las credenciales de la api & El sistema muestra en el mapa la ubicación de los 42 camiones propios, igual que hoy, pero integrado en la nueva solución. & Crítica / E1 &  \\
  RF-14.02 & Gerente de Operaciones / Planificador de Rutas & El sistema debe mostrar en el mapa las rutas planificadas vs las rutas reales recorridas por cada camión, para identificar desviaciones y optimizar las rutas futuras. & El sistema contiene las rutas planificadas para el día & El gerente de operaciones ve que el camión 5 se desvía 5 km en su ruta; investiga si es por tráfico, error del conductor o una parada no registrada. & Alta / E1 &  \\
  RF-14.03 & Gerente de Operaciones & El sistema debe generar una alerta cuando el camión se desvíe más de 3 km de la ruta planificada, para que el centro de control pueda reaccionar. & La ruta planificada del camión está cargada en el sistema & El centro de control recibe una alerta si un camión se desvía para investigar si hay un problema (accidente, robo, etc.). & Alta / E1 &  \\
  RF-14.04 & Gerente de Finanzas / Gerente de Operaciones & El sistema debe integrar los datos de telemetría con el módulo de costo de servir, imputando kilómetros recorridos y tiempo de viaje a cada entrega y cliente. & Los datos de telemetría pueden vincularse a entregas específicas y clientes. & El costo de servir por cliente incluye el costo real de los kilómetros recorridos para llegar a su local. & Crítica / E2 &  \\
  RF-14.05 & Tripulación propia y peoneta / Conductores externos / Empresas Transp. & El sistema debe registrar qué conductor propio o de transportista externo está asociado a cada camión en cada viaje, consumiendo la identidad validada por el mecanismo de autenticación (RF-15.02). Esta vinculación es necesaria para imputar responsabilidades en cobros, entregas y eventos de ruta, y para resolver la Decisión \#5. & El conductor ha sido autenticado correctamente (RF-06.08 / RF-15.02), el vehículo está asignado a un viaje en el sistema. & Se resuelve la Decisión \#5; el sistema exige que el conductor se identifique (PIN o QR) antes de iniciar la ruta, incluso si es externo. & Alta / E1 & ), Cap. 4.6 \\
  RF-14.06 & Gerente de Operaciones / Tripulación propia y peoneta / Conductores externos & El sistema debe permitir la gestión de geocercas zonas geográficas, para identificar cuándo un camión entra o sale de una zona ej: CD Talca, zona de clientes, para medir tiempos de llegada y salida. & El usuario tiene permisos para crear y administrar geocercas, el sistema posee un mapa base o capacidad de definir polígonos sobre coordenadas & El sistema registra automáticamente la hora de llegada del camión al área del cliente, para validar el cumplimiento de la ventana horaria. & Media / E1 &  \\
\end{longtable}}
